\chapter{Semântica operacional para estado}

\section{Introdução}

No capítulo anterior, estudamos a semântica operacional de uma
linguagem de expressões aritméticas puras. Naquela linguagem, o valor
de uma expressão depende apenas da sua estrutura sintática: não há
variáveis, não há memória, não há efeitos colaterais. A execução é,
portanto, uma função pura: dada a expressão, o resultado é unicamente
determinado.

Programas reais, contudo, manipulam \textbf{estado}: valores são
armazenados em variáveis, lidos da entrada e escritos na saída. Para
modelar formalmente esse comportamento, a semântica operacional
precisa ser estendida com um componente que representa o estado da
memória em cada momento da execução.

Este capítulo estende a linguagem de expressões com
\textbf{variáveis} e apresenta uma linguagem de \textbf{comandos}
simples, a linguagem \emph{Line},  que inclui atribuição,
leitura da entrada padrão e impressão. Definimos a semântica
small-step e big-step para essa linguagem e mostramos como a
especificação semântica se traduz diretamente no interpretador
disponível em
\begin{flushleft}
  \verb|Line/Interp/LineInterp.hs|.
\end{flushleft}

\section{A Linguagem Line}

A linguagem \emph{Line} é uma linguagem imperativa mínima. Um
programa é uma sequência de comandos (instruções) que são executados
em ordem. A sua sintaxe abstrata é:

\[
  \begin{array}{lcl}
    \mathit{prog} & ::= & s^*\\
    s & ::= & x \mathbin{:=} e \\
    & \mid & \mathbf{read}\; x \\
    & \mid & \mathbf{print}\; e\\
    e & ::= & n \mid x \mid e_1 + e_2 \mid e_1 * e_2
  \end{array}
\]

onde \(n\) é um inteiro, \(x\) é um nome de variável e \(s^*\) denota
uma sequência (possivelmente vazia) de comandos.

Em Haskell, esses tipos são definidos como:

\begin{minted}{haskell}
data Line = Line [Stmt]

data Stmt
  = SAssign Var Exp
  | SRead   Var
  | SPrint  Exp

data Exp
  = EInt Int
  | EVar Var
  | Exp :+: Exp
  | Exp :*: Exp

type Var = String
\end{minted}

\section{Ambientes}

Para dar significado a um programa que usa variáveis, precisamos de
uma noção de \textbf{ambiente} (uma representação da memória): uma
função parcial \(\sigma : \mathit{Var}
\rightharpoonup \mathbb{Z}\) que mapeia nomes de variáveis para os
seus valores inteiros atuais.

Escrevemos:

\begin{itemize}
  \item \(\sigma(x)\) para o valor de \(x\) em \(\sigma\) (supondo \(x
    \in \mathrm{dom}(\sigma)\))
  \item \(\sigma[x \mapsto n]\) para o ambiente obtido de \(\sigma\)
    atualizando (ou inserindo) a variável \(x\) com o valor \(n\)
  \item \(\emptyset\) para o ambiente vazio (sem variáveis definidas)
\end{itemize}

Em Haskell, o ambiente é representado como
\haskell{Map String Int} do módulo
\haskell{Data.Map}.

\section{Semântica Small-Step}

\subsection{Expressões}

A avaliação de expressões não modifica o ambiente: expressões são
\textbf{puras} no sentido de que apenas consultam \(\sigma\), sem
alterá-lo. Uma configuração de expressão é um par \((\sigma, e)\) e a
relação de redução é \((\sigma, e) \to (\sigma, e')\).

\textbf{Variável:}

\[
  \infer[_{\{Var\}}]{(\sigma,\, x) \to (\sigma,\, \sigma(x))}{x \in
  \mathrm{dom}(\sigma)}
\]

Uma variável é substituída pelo seu valor corrente no ambiente. Se
\(x\) não estiver definida em \(\sigma\), a avaliação trava: erro
de variável não inicializada.

\textbf{Redução de subexpressões:}

\[
  \begin{array}{cc}
    \infer[_{\{AddL\}}]{(\sigma,\, e_1 + e_2) \to (\sigma,\, e_1' +
    e_2)}{(\sigma, e_1) \to (\sigma, e_1')}
    &
    \infer[_{\{AddR\}}]{(\sigma,\, n_1 + e_2) \to (\sigma,\, n_1 +
    e_2')}{(\sigma, e_2) \to (\sigma, e_2')}
  \end{array}
\]

\[
  \infer[_{\{Add\}}]{(\sigma,\, n_1 + n_2) \to (\sigma,\, n_1
  \mathbin{\hat{+}} n_2)}{}
\]

As regras para multiplicação são análogas: (MulL), (MulR) e (Mul).

\subsection{Comandos}

A configuração de um comando inclui o ambiente, pois os comandos
podem modificá-lo. Usamos a notação \((\sigma, s) \to (\sigma',
s')\), onde \(s'\) é o comando residual após a redução. Quando um
comando termina sua execução, dizemos que ele reduz para
\(\mathbf{skip}\), um comando que não faz nada.

\textbf{Atribuição: avalia a expressão primeiro:}

\[
  \infer[_{\{AssignStep\}}]{(\sigma,\, x \mathbin{:=} e)\to
  (\sigma,\, x \mathbin{:=} e')}{(\sigma, e) \to (\sigma, e')}
\]

\textbf{Atribuição: expressão já é um valor:}

\[
  \infer[_{\{Assign\}}]{(\sigma,\, x \mathbin{:=} n) \to (\sigma[x \mapsto n],\,
  \mathbf{skip})}{}
\]

\textbf{Impressão: avalia a expressão primeiro:}

\[
  \infer[_{\{PrintStep\}}]{(\sigma,\, \mathbf{print}\; e)\to
  (\sigma,\, \mathbf{print}\; e')}{(\sigma, e) \to (\sigma, e')}
\]

\textbf{Impressão: expressão já é um valor:}

\[
  \infer[_{\{Print\}}]{(\sigma,\, \mathbf{print}\; n) \to (\sigma,\,
  \mathbf{skip})}{}
  \quad [\text{ emite $n$ na saída}]
\]

\textbf{Sequência:}

\[
  \infer[_{\{SeqSkip\}}]{(\sigma,\, \mathbf{skip};\, s^\star) \to
  (\sigma,\, s^\star)}{}
\]

\[
  \infer[_{\{Seq\}}]{(\sigma,\, s;\, s^\star) \to (\sigma',\, s';\,
  s^\star)}{(\sigma, s) \to (\sigma', s')}
\]

\subsection{Exemplo: Small-Step}

Considere o programa:

\begin{verbatim}
x := 2 + 3
print x * 4
\end{verbatim}

Iniciamos com o ambiente vazio \(\emptyset\). A execução procede assim:

\[
  (\emptyset,\; x \mathbin{:=} 2 + 3;\; \mathbf{print}\; x * 4)
\] \[
  \xrightarrow{\text{Seq, Add}} ([x \mapsto 5],\; \mathbf{skip};\;
  \mathbf{print}\; x * 4)
\] \[
  \xrightarrow{\text{SeqSkip}} ([x \mapsto 5],\; \mathbf{print}\; x * 4)
\] \[
  \xrightarrow{\text{PrintStep, Var}} ([x \mapsto 5],\; \mathbf{print}\; 5 * 4)
\] \[
  \xrightarrow{\text{PrintStep, Mul}} ([x \mapsto 5],\; \mathbf{print}\; 20)
\] \[
  \xrightarrow{\text{Print}} ([x \mapsto 5],\; \mathbf{skip}) \quad
  \text{[emite 20]}
\]

Cada passo é justificado por uma regra; a sequência termina quando
restam apenas \(\mathbf{skip}\)s.

\section{Semântica Big-Step}

Na semântica big-step, abstraímos os passos intermediários e
definimos diretamente a relação entre configuração inicial e resultado final.

\subsection{Expressões}

A relação \(\sigma \vdash e \Downarrow n\) significa ``no ambiente
\(\sigma\), a expressão \(e\) avalia para o inteiro \(n\)''.

\textbf{Número:}

\[
  \infer[_{\{Num\}}]{\sigma \vdash n \Downarrow n}{}
\]

\textbf{Variável:}

\[
  \infer[_{\{Var\}}]{\sigma \vdash x \Downarrow \sigma(x)}{x \in
  \mathrm{dom}(\sigma)}
\]

\textbf{Adição:}

\[
  \infer[_{\{Add\}}]{\sigma \vdash e_1 + e_2 \Downarrow n_1
  \mathbin{\hat{+}} n_2}{\sigma \vdash e_1 \Downarrow n_1 & \sigma
  \vdash e_2 \Downarrow n_2}
\]

\textbf{Multiplicação:}

\[
  \infer[_{\{Mul\}}]{\sigma \vdash e_1 * e_2 \Downarrow n_1
  \mathbin{\hat{*}} n_2}{\sigma \vdash e_1 \Downarrow n_1 & \sigma
  \vdash e_2 \Downarrow n_2}
\]

\subsection{Comandos}

A relação \(\sigma, s \Downarrow \sigma'\) significa ``executar o
comando \(s\) no ambiente \(\sigma\) produz o ambiente \(\sigma'\)''.
O ambiente é o único ``resultado'' de um comando; os efeitos de I/O
(impressão, leitura) são considerados efeitos colaterais implícitos.

\textbf{Atribuição:}

\[
  \infer[_{\{Assign\}}]{\sigma,\; x \mathbin{:=} e \;\Downarrow\;
  \sigma[x \mapsto n]}{\sigma \vdash e \Downarrow n}
\]

\textbf{Impressão:}

\[
  \infer[_{\{Print\}}]{\sigma,\; \mathbf{print}\; e \;\Downarrow\;
  \sigma}{\sigma \vdash e \Downarrow n}
\]

O ambiente não muda com \haskell{print}; apenas o
valor de \(e\) é emitido como efeito colateral.

\textbf{Leitura:}

\[
  \infer[_{\{Read\}}]{\sigma,\; \mathbf{read}\; x \;\Downarrow\;
  \sigma[x \mapsto n]}{}
  \quad n \text{ lido da entrada}
\]

A regra (Read) é um axioma cuja premissa implícita é a leitura de um
inteiro \(n\) da entrada padrão.

\textbf{Sequência vazia:}

\[
  \infer[_{\{SeqNil\}}]{\sigma,\; \epsilon \;\Downarrow\; \sigma}{}
\]

\textbf{Sequência não-vazia:}

\[
  \infer[_{\{SeqCons\}}]{\sigma,\; s;\, s^* \;\Downarrow\;
  \sigma''}{\sigma,\; s \;\Downarrow\; \sigma' \quad \sigma',\; s^*
  \;\Downarrow\; \sigma''}
\]

\subsection{Exemplos: Big-Step}\label{exemplos-big-step}

\textbf{Exemplo 1:} Derivação para \haskell{x := 2 +
3} partindo de \(\emptyset\).

\[
  \dfrac{
    \dfrac{
      \dfrac{}{\emptyset \vdash 2 \Downarrow 2}
      \quad
      \dfrac{}{\emptyset \vdash 3 \Downarrow 3}
    }{\emptyset \vdash 2 + 3 \Downarrow 5}
  }{\emptyset,\; x \mathbin{:=} 2 + 3 \;\Downarrow\; [x \mapsto 5]}
\]

\textbf{Exemplo 2:} Derivação para o programa completo
(\haskell{x := 2 + 3; print x * 4}) partindo de \(\emptyset\).

\[
  \dfrac{
    \underbrace{\emptyset,\; x \mathbin{:=} 2+3 \;\Downarrow\; [x
    \mapsto 5]}_{\text{Exemplo 1}}
    \quad
    \dfrac{
      \dfrac{
        \dfrac{}{[x\mapsto 5] \vdash x \Downarrow 5}
        \quad
        \dfrac{}{[x\mapsto 5] \vdash 4 \Downarrow 4}
      }{[x\mapsto 5] \vdash x \star 4 \Downarrow 20}
    }{[x\mapsto 5],\; \mathbf{print}\; x\star 4 \;\Downarrow\; [x\mapsto 5]}
  }{\emptyset,\; x \mathbin{:=} 2+3;\; \mathbf{print}\; x\star 4
  \;\Downarrow\; [x\mapsto 5]}
\]

A derivação confirma que, ao final, \(x = 5\) e o valor \(20\) foi emitido.

\section{Implementação em Haskell}

\subsection{O Ambiente como Mônada de Estado}

A semântica big-step dita que cada comando recebe um ambiente e
produz um (possivelmente diferente) ambiente. Em Haskell, isso é
capturado pela \textbf{mônada de estado} (\haskell{StateT}):

\begin{minted}{haskell}
type Env = Map String Int
type Interp a = ExceptT String (StateT Env IO) a
\end{minted}

A pilha de mônadas combina três capacidades:
\begin{itemize}
  \item \haskell{StateT Env}: o ambiente \(\sigma\)
    implícito, que pode ser lido com \haskell{get} e
    atualizado com \haskell{modify};
  \item  \haskell{ExceptT String}: tratamento de erros
    (como variável não definida) via \haskell{throwError}
  \item \haskell{IO}: efeitos de I/O para
    \haskell{read} e \haskell{print}.
\end{itemize}

O ponto de entrada executa a pilha a partir do ambiente vazio:

\begin{minted}{haskell}
interp :: Line -> IO ()
interp p = do
    (r, _) <- runStateT (runExceptT (interpProgram p)) Map.empty
    case r of
      Left err -> putStrLn err
      Right _  -> pure ()
\end{minted}

\subsection{Interpretador de Expressões}

A função \haskell{interpExp} implementa as regras
(Num), (Var), (Add) e (Mul):

\begin{minted}{haskell}
interpExp :: Exp -> Interp Int
interpExp (EInt n) = pure n
interpExp (EVar v) = askVar v
interpExp (e1 :+: e2) = (+) <$> interpExp e1 <*> interpExp e2
interpExp (e1 :*: e2) = (*) <$> interpExp e1 <*> interpExp e2
\end{minted}

A função auxiliar \haskell{askVar} implementa a regra
(Var): consulta o ambiente corrente e lança um erro se a variável não
estiver definida (caso não coberto pelas regras semânticas):

\begin{minted}{haskell}
askVar :: String -> Interp Int
askVar s = do
    env <- get
    case Map.lookup s env of
      Nothing  -> throwError $ unwords ["Undefined variable:", s]
      Just val -> pure val
\end{minted}

\subsection{Interpretador de Comandos}

A função \haskell{interpStmt} implementa as regras
(Assign), (Print) e (Read):

\begin{minted}{haskell}
interpStmt :: Stmt -> Interp ()
interpStmt (SAssign v e) = do
    val <- interpExp e
    updateVar v val
interpStmt (SRead v) = do
    str <- liftIO getLine
    let val = read str
    updateVar v val
interpStmt (SPrint e) = do
    val <- interpExp e
    liftIO $ print val
\end{minted}

A correspondência com as regras semânticas é direta:

\begin{itemize}
  \item \haskell{SAssign v e}: avalia
    \haskell{e} (regra Assign, premissa) e
    atualiza o ambiente com \haskell{updateVar}
    (conclusão: \(\sigma[x \mapsto n]\)).
  \item \haskell{SPrint e}: avalia
    \haskell{e} e imprime o resultado; o ambiente não
    muda (regra Print).
  \item \haskell{SRead v}: lê um inteiro da entrada
    padrão e armazena em \haskell{v} (regra Read).
\end{itemize}

A função \haskell{updateVar} encapsula a operação
\(\sigma[x \mapsto n]\):

\begin{minted}{haskell}
updateVar :: String -> Int -> Interp ()
updateVar s val = modify (Map.insert s val)
\end{minted}

A interpretação de um programa aplica
\haskell{interpStmt} a cada comando em sequência,
correspondendo às regras (SeqNil) e (SeqCons):

\begin{minted}{haskell}
interpProgram :: Line -> Interp ()
interpProgram (Line ss) = mapM_ interpStmt ss
\end{minted}

\haskell{mapM_} sequencia os comandos na ordem e
descarta os resultados unitários, o que equivale exatamente a
encadear as regras (SeqCons) até atingir (SeqNil).

\subsection{Executando o Interpretador}

\begin{verbatim}
cabal run line -- --interp --file programa.line
\end{verbatim}

Por exemplo, dado o arquivo \haskell{programa.line}:

\begin{verbatim}
x := 10
y := x * 2
print y
\end{verbatim}

O interpretador produziria 20.

\section{Conclusão}

Neste capítulo, estendemos a semântica operacional para lidar com
\textbf{estado}, introduzindo ambientes que mapeiam variáveis a
valores. As contribuições principais foram:

\begin{itemize}
  \item
    A \textbf{semântica small-step} para comandos usa configurações
    \((\sigma, s)\) que transportam o ambiente ao longo dos passos de
    redução. A regra (Assign) modifica o ambiente na transição; as
    outras regras o propagam sem alteração.
  \item
    A \textbf{semântica big-step} relaciona um comando e um ambiente
    inicial ao ambiente final, abstraindo os passos intermediários. A
    sua estrutura de ``entrada → saída de ambiente'' traduz-se
    diretamente na mônada de estado de Haskell.
  \item
    A \textbf{implementação} em
    LineInterp.hs usa a pilha
    \haskell{ExceptT String (StateT Env IO)}, que
    integra estado (ambiente), erros e I/O. Cada equação do
    interpretador corresponde a uma regra semântica.
\end{itemize}

O padrão estado-como-mônada que vimos aqui reaparecerá nos próximos
capítulos quando estendermos a linguagem com estruturas de controlo
como condicionais e laços.
